%% --- Project Summary (1 page) ---
%% NSF requires three labeled subsections: Overview, Intellectual Merit, Broader Impacts.

\begin{center}
  {\large\bfseries Project Summary}\\[0.3em]
  {\large\bfseries \sysname: AI-Enabled Digital Twins for Continuous Red Teaming of Open-Source Critical-Infrastructure Stacks}
\end{center}

\noindent\textbf{Overview.}
Modern critical-infrastructure services---hospital records, financial
exchanges, telecommunications back-ends, government portals, cloud
platforms, utility management systems---increasingly run on complex
stacks of open-source software: databases, message brokers, container
orchestrators, service meshes, identity providers, web servers, and
custom glue code. Yet assessing the security posture of a specific
\emph{deployment} of such a stack remains a largely manual, ad hoc, and
infrequent activity.
Open-source infrastructure stacks are inherently difficult to analyze
in situ: they are heterogeneous assemblies of dozens of components,
each with its own configuration surface and version history; they are
dynamic in both composition and usage; and they execute services whose
availability must not be disrupted. An aggressive vulnerability scan
that is routine in a lab environment can, in production, trigger
cascading failures in a database cluster, corrupt a message queue, or
violate a regulated uptime SLA. As a result, security assessments are
scheduled rarely and performed by scarce human experts, while
adversaries---now increasingly augmented with AI---operate continuously.

This project proposes \sysname, an agentic framework that enables the
\emph{continuous, in-depth, and non-disruptive} security assessment of
open-source infrastructure stacks. \sysname uses cooperating teams of
AI agents to
(i)~ingest documentation, configuration, and runtime signals from a
target deployment into a structured and vectorized knowledge base;
(ii)~passively observe the system to build an evolving model of its
components, software, and flows; (iii)~synthesize a \emph{digital twin}
that is a minimally representative yet faithful replica of the original;
(iv)~continuously red-team the twin using component-level analyses
(fuzzing, static analysis, symbolic execution) composed into multi-step
attack plans; and (v)~surface prioritized findings to human analysts,
who validate them on the original system and whose feedback refines
the twin. Four agent roles---Knowledge Base, Observer, Replication,
and Red Team---operate as a feedback loop rather than a linear
pipeline, so that twin and analysis co-evolve with the real system.

\smallskip
\noindent\textbf{Intellectual Merit.}
\sysname advances the state of the art in three ways. First, the project
develops new techniques for automatically constructing a living model of
an open-source infrastructure deployment from heterogeneous, partial,
and sometimes inconsistent evidence, combining graph-structured
knowledge with natural-language representations to support agent
reasoning. Second, it formulates digital
twin synthesis as a \emph{controlled-fidelity replication problem}, in
which diversity of component classes, topological patterns, and data
flows are explicit quantifiable parameters; this makes it feasible to
replicate large deployments at a useful level of abstraction rather
than requiring a bit-for-bit copy of every instance and every dataset.
Third, the project builds on
the PIs' prior work on opaque-system modeling (Conware, HALucinator), AI
planning for privilege escalation (ChainReactor), and autonomous
vulnerability discovery (Artiphishell, developed for the DARPA AIxCC
competition) to compose single-component vulnerabilities into end-to-end
multi-step attacks whose plans are expressed in a structured planning
language and executed against the twin without risk to the real system.
Together, these contributions create a new scientific foundation for
continuous, automated, adversarial assessment of the open-source
software that runs critical infrastructure.

\smallskip
\noindent\textbf{Broader Impacts.}
The security of open-source critical-infrastructure stacks is a matter
of direct national importance: attacks on healthcare systems,
financial services, energy operators, and government platforms have
already caused measurable harm to U.S.\ communities, and
the gap between attacker capability and defender capacity continues to
widen. By operating continuously and without disrupting operations,
\sysname can substantially raise the baseline resilience of critical
services whose operators cannot currently afford frequent traditional
red-team engagements. The project will release its framework as open
source, following the PIs' established practice with widely-used
artifacts such as \texttt{angr}. The PIs will integrate the research
into UCSB graduate and undergraduate curricula, involve high-school,
REU, and ERSP students (including students from underrepresented
groups), and organize an annual CTF competition, focused on the
security of open-source infrastructure, open to
the public. Industry collaboration, internships, and publication in
top-tier security venues will ensure that results transfer rapidly to
practitioners defending real systems.

\smallskip
\noindent\textbf{Keywords:} open-source software; critical
infrastructure; digital twins; AI agents; red teaming; continuous
security assessment.
